\section[Self-Supervised Learning]{Self-Supervised Learning as Predictive Representation Learning}

Section 4.1 introduced the chapter's main thesis: modern machine learning increasingly seeks a representation
\[
\phi : X \to Z
\]
that can be reused across many downstream tasks.
The next question is where such a representation comes from when task-specific labels are scarce.
One major answer is \emph{self-supervised learning}.
In self-supervision, the learning signal is extracted from structure already present in the data.
The model is not given an external class label for every example.
Instead it is trained to solve a predictive problem constructed from the example itself.

\medskip

This section is one of the conceptual cores of the chapter because it explains why modern pretraining is not merely ``training without labels.''
It is better understood as \emph{predictive representation learning}.
The pretraining objective is chosen so that success requires the model to organize inputs into internal coordinates that preserve reusable regularities.
In language, those regularities may include syntax, discourse, and semantic compatibility.
In vision, they may include shape, object identity, spatial composition, and invariances to nuisance variation.
In temporal data, they may include latent state and dynamical structure.

\subsection*{An abstract view of self-supervised learning}

Let $X$ be an input space and let $P_X$ be a distribution over raw observations.
A self-supervised system first constructs from $x \sim P_X$ one or more derived objects: masked views, augmented views, contexts, future targets, or paired modalities.
These derived objects define a prediction problem.
The representation map $\phi_\theta$ and possibly an auxiliary head $h_\theta$ are then trained to minimize a loss that depends only on these internally constructed targets.

\begin{definition}[Abstract self-supervised objective]
A self-supervised learning problem is an optimization problem of the form
\[
\min_{\theta}
\; \mathbb{E}_{x \sim P_X}
\; \mathbb{E}_{\omega \sim Q(\cdot \mid x)}
\big[
\ell_{\mathrm{ssl}}\big(\phi_\theta, h_\theta; x, \omega\big)
\big],
\]
where $\omega$ is an auxiliary random object derived from $x$ and sampled from a rule $Q(\cdot \mid x)$.
The variable $\omega$ may encode a mask, an augmentation, a context window, a time index, or a pairing relation.
The loss is self-supervised when its target is determined from $x$ and $\omega$ rather than from an externally supplied task label.
\end{definition}

This abstraction is useful because it separates the enduring mathematical structure from implementation details.
Every self-supervised method chooses:
\begin{itemize}
    \item what predictive relation should be extracted from the data,
    \item what family of transformations or corruptions defines the task,
    \item what representation must be learned to solve that task.
\end{itemize}
In this sense self-supervision is not one method but a design pattern.

\subsection*{A taxonomy of self-supervised objectives}

Three large families organize most modern self-supervised learning: masked prediction, autoregressive prediction, and contrastive learning.
They differ in what is hidden, what is predicted, and what kind of invariance or dependency structure is encouraged.

\paragraph{1. Masked prediction.}
Part of the input is hidden and the model is asked to infer the missing content from the visible part.
If $x = (x_1,\dots,x_T)$ is a structured input and $M \subseteq \{1,\dots,T\}$ is a set of masked positions, we write the visible context as $x_{\bar M}$ and the masked content as $x_M$.
A canonical masked objective is
\[
\mathcal{L}_{\mathrm{mask}}(\theta)
=
\mathbb{E}_{x,M}
\big[
-\log p_\theta(x_M \mid x_{\bar M}, M)
\big].
\]
This trains the model to reconstruct hidden information from context.
The objective does not directly tell the model what semantic abstractions to learn, but it does force the representation to capture dependencies among observed parts.

\paragraph{2. Autoregressive prediction.}
In sequential data one may ask the model to predict each token from the past.
For a sequence $x_1,\dots,x_T$, the objective is
\[
\mathcal{L}_{\mathrm{AR}}(\theta)
=
\mathbb{E}_{x}
\Big[
\sum_{t=1}^{T}
-\log p_\theta(x_t \mid x_{<t})
\Big].
\]
This is a standard conditional log-likelihood.
It encourages the representation at time $t$ to summarize the preceding context in a way useful for continuation.
Next-token prediction in language models is the most visible modern example.

\paragraph{3. Contrastive learning.}
Rather than reconstructing hidden content, contrastive learning distinguishes compatible pairs from incompatible ones.
Let $a$ and $a'$ be random augmentations, and let
\[
u = a(x), \qquad v = a'(x)
\]
be two views of the same underlying example.
Let $\tilde v_1,\dots,\tilde v_K$ be views of other examples.
With an encoder $\phi_\theta$ and similarity function $s(\cdot,\cdot)$, a typical contrastive loss takes the form
\[
\mathcal{L}_{\mathrm{con}}(\theta)
=
\mathbb{E}
\left[
-\log
\frac{\exp(s(z,z^+)/\tau)}
{\exp(s(z,z^+)/\tau)+\sum_{k=1}^{K} \exp(s(z,z_k^-)/\tau)}
\right],
\]
where $z=\phi_\theta(u)$, $z^+=\phi_\theta(v)$, and $z_k^- = \phi_\theta(\tilde v_k)$.
The temperature $\tau>0$ controls how sharply similarities are weighted.
The objective pushes together representations of related views and pushes apart representations of unrelated examples.

\subsection*{Positives, negatives, augmentations, and contexts}

The notation above is standard enough that it deserves explicit definitions.

\begin{definition}[Augmentation]
An \emph{augmentation} is a transformation $a : X \to X'$ that preserves enough of the underlying content for the transformed object to remain a valid view of the same example for the pretraining task.
In images, augmentations may include cropping, color perturbation, or flipping.
In language, they may be more constrained because arbitrary edits can destroy meaning.
\end{definition}

\begin{definition}[Positive pair and negative pair]
In contrastive learning, a \emph{positive pair} is a pair of views $(u,v)$ deemed compatible because they are generated from the same underlying source example or matched context.
A \emph{negative pair} is a pair $(u,\tilde v)$ deemed incompatible because the two views come from different source examples or mismatched contexts.
\end{definition}

\begin{definition}[Context]
A \emph{context} is the observed part of a structured input that is used to predict another part.
For masked learning the context is typically $x_{\bar M}$.
For autoregression the context at position $t$ is the prefix $x_{<t}$.
More generally, a context is any information retained after a task-specific corruption or restriction has been applied.
\end{definition}

These definitions matter because a large part of self-supervised design is hidden inside them.
Choosing an augmentation implicitly decides what information the representation should ignore.
Choosing positives and negatives decides what counts as semantic compatibility.
Choosing a context rule decides what dependencies the model is forced to model.

\subsection*{What contrastive learning is trying to do}

There are two complementary ways to understand contrastive learning.
One is geometric: views of the same object should map nearby, while views of different objects should not collapse into one indistinguishable cluster.
The other is information-theoretic: the representation of one view should retain information predictive of another compatible view.
A standard formal statement linking these viewpoints is the InfoNCE lower-bound result.

\begin{proposition}[InfoNCE lower-bound intuition]
Let $(U,V)$ be jointly distributed random variables, and let $V_1,\dots,V_N$ be candidates such that one candidate is paired with $U$ according to the joint law and the others are sampled independently from the marginal law of $V$.
If a scoring function $f_\theta(U,V)$ is trained with the multiclass contrastive objective
\[
\mathcal{L}_{\mathrm{NCE}}(\theta)
=
\mathbb{E}
\left[
-\log
\frac{\exp(f_\theta(U,V_1))}
{\sum_{j=1}^{N} \exp(f_\theta(U,V_j))}
\right],
\]
then, under the standard sampling construction above,
\[
I(U;V) \ge \log N - \mathcal{L}_{\mathrm{NCE}}(\theta)
\]
for the optimal scoring rule and, more generally, low contrastive loss is evidence that the score retains information useful for identifying the correct paired view among negatives.
\end{proposition}

\paragraph{Proof sketch and interpretation.}
The classification problem hidden inside InfoNCE asks the model to identify which candidate $V_j$ is paired with $U$.
If $U$ and $V$ share substantial mutual information, then this identification problem is easier than chance.
The optimal logit is related to a log-density ratio between the joint distribution and the product of marginals.
Substituting that optimal score into the cross-entropy objective yields the inequality above.
The proposition should be read carefully.
It does not say that every practical contrastive pipeline estimates mutual information well.
It says that under a mathematically clean negative-sampling construction, successful contrastive discrimination is linked to preserving shared information between two views.

This proposition is useful because it elevates contrastive learning above pure heuristic.
There is a real theorem-level bridge from a discrimination objective to an information quantity.
At the same time, the theorem does not automatically settle practice.
Finite batch size, imperfect negatives, architecture choice, projector heads, and augmentation design all affect what the learned representation actually becomes.

\subsection*{What is known, what is empirical, and what remains heuristic}

Because self-supervised learning is now a central methodology, it is important to distinguish three layers of understanding.

\paragraph{What is mathematically well grounded.}
Masked prediction and autoregressive prediction are conditional likelihood objectives on induced prediction tasks.
The optimization problem is explicit.
If the model class is rich enough and optimization succeeds, the optimum is tied to modeling a conditional distribution.
Contrastive learning also has formal analyses, including the lower-bound interpretation above and geometric analyses of alignment and non-collapse in simplified settings.
Thus self-supervision is not merely folklore.

\paragraph{What is strongly supported empirically.}
Large-scale pretraining can produce representations that transfer remarkably well across downstream tasks.
Masked language modeling, next-token prediction, and contrastive image learning have all been shown in practice to learn features useful far beyond the pretraining objective itself.
Linear probing and fine-tuning experiments repeatedly demonstrate that pretrained representations encode substantial reusable structure.
These are empirical facts about current systems.

\paragraph{What remains design practice or heuristic.}
Many decisive choices still come from engineering judgment rather than closed theory.
How much masking should be used?
What augmentations preserve semantic identity rather than destroy it?
How many negatives are helpful?
What temperature should be used in contrastive loss?
Why do predictor heads and stop-gradient tricks help in some image methods?
Why do some objectives transfer better than others even when training loss looks similar?
These questions are only partially understood.
The current field combines theorem, experiment, and craft.

\subsection*{Worked example: masked language modeling}

A compact toy example makes the predictive structure concrete.
Consider the sentence
\[
\text{``the cat sat on the mat''}.
\]
Suppose positions $3$ and $6$ are masked.
The visible context is therefore
\[
(\text{the},\text{cat},\texttt{[MASK]},\text{on},\text{the},\texttt{[MASK]}).
\]
A masked language model produces contextual representations
\[
h_t = \phi_\theta(x_{\bar M},M)_t
\]
for each position $t$, and then predicts the token at each masked position via a softmax head:
\[
p_\theta(x_t = w \mid x_{\bar M},M)
=
\frac{\exp(W_w^{\top} h_t + b_w)}{\sum_{w'} \exp(W_{w'}^{\top} h_t + b_{w'})}.
\]
If the correct targets are ``sat'' and ``mat,'' the contribution to the loss is
\[
-\log p_\theta(\text{sat} \mid x_{\bar M},M, t=3)
-\log p_\theta(\text{mat} \mid x_{\bar M},M, t=6).
\]

Why can this teach useful representation rather than mere word filling?
Because correct prediction requires the model to exploit several kinds of dependency at once.
At position $3$, the local context ``the cat \_ on'' strongly suggests a verb compatible with the subject.
At position $6$, the phrase ``on the \_'' suggests a noun, but selecting ``mat'' over alternatives is easier if the model also preserves coherence with the earlier words.
Thus the internal state must encode grammatical and semantic compatibility across the sentence.

\begin{example}[What is learned in the toy masked example]
The model is not explicitly told to learn syntax, parts of speech, or semantic role.
Those structures emerge as useful internal variables because they help predict the masked tokens.
This is a recurring pattern in self-supervision: the objective names one task, but the representation internalizes latent regularities needed to solve it.
\end{example}

One should also note a limitation.
Masked prediction is not the same as direct generation.
It teaches the model to use bidirectional context for infilling, not necessarily to produce a left-to-right continuation policy.
This is one reason why different self-supervised objectives can lead to different downstream strengths.

\subsection*{Worked example: contrastive image learning}

Now consider a toy image example.
Let $x$ be an image of a bicycle.
Apply two random augmentations:
\[
u = a(x), \qquad v = a'(x),
\]
where one augmentation crops the left side of the bicycle and the other changes color and brightness.
Although the two resulting views are visually different, they still correspond to the same underlying object.
Take another image $x'$ of a dog and form a negative view $\tilde v = a''(x')$.
Suppose the encoder produces normalized embeddings
\[
z = \phi_\theta(u),
\qquad
z^+ = \phi_\theta(v),
\qquad
z^- = \phi_\theta(\tilde v),
\]
and use dot-product similarity $s(z,z') = z^{\top}z'$.

Assume, for illustration, that
\[
s(z,z^+) = 2.0,
\qquad
s(z,z^-) = 0.5,
\qquad
\tau = 1.
\]
With one negative, the contrastive loss for this anchor is
\[
-\log
\frac{e^{2.0}}{e^{2.0}+e^{0.5}}
\approx 0.20.
\]
If the positive similarity were smaller, say $0.8$, then the loss would be
\[
-\log
\frac{e^{0.8}}{e^{0.8}+e^{0.5}}
\approx 0.55,
\]
which is worse.
Thus the training rule explicitly rewards making the two bicycle views more compatible than the bicycle--dog pair.

The important conceptual point is that the representation is being shaped by invariance assumptions hidden in the augmentations.
If cropping and color perturbation preserve object identity, then a good representation should not change radically under those transformations.
But this is also where design enters.
An augmentation that destroys semantic content teaches the wrong invariance.
For example, an extreme crop that removes the bicycle entirely should not remain a positive view in any meaningful sense.

\begin{figure}[h]
\centering
\fbox{%
\begin{minipage}{0.94\textwidth}
\centering
\textbf{Two common self-supervised pipelines}

\vspace{0.5em}
\begin{tabular}{p{0.43\textwidth}p{0.43\textwidth}}
\textbf{Masked prediction} & \textbf{Contrastive learning} \\
Raw example $x$ & Raw example $x$ \\
$\downarrow$ apply mask $M$ & $\downarrow$ sample augmentations $a,a'$ \\
Visible context $x_{\bar M}$ & Two views $u=a(x),\; v=a'(x)$ \\
$\downarrow$ encoder $\phi_\theta$ & $\downarrow$ shared encoder $\phi_\theta$ \\
Contextual states $h_t$ & Embeddings $z, z^+$ \\
$\downarrow$ predict hidden content $x_M$ & $\downarrow$ compare with negatives $z_1^-,\dots,z_K^-$ \\
Loss: reconstruction / conditional likelihood & Loss: align positives, separate negatives
\end{tabular}
\end{minipage}}
\caption{A monochrome comparison of two major self-supervised pipelines.
Masked methods learn by recovering omitted information from context.
Contrastive methods learn by making compatible views more similar than incompatible ones.}
\end{figure}

\subsection*{Why self-supervision can support transfer}

We can now state the chapter-level intuition more sharply.
Self-supervised learning tends to help downstream tasks when solving the pretext objective requires the model to preserve regularities that many downstream tasks also need.
The mechanism is not magical.
A representation becomes reusable when the pretraining objective rewards internal coordinates that expose broadly useful structure.

Masked prediction often encourages context-sensitive features.
Autoregression encourages representations that summarize history and continuation constraints.
Contrastive learning often encourages invariances to nuisance transformations while preserving semantic identity.
These are not identical biases.
Each objective selects a somewhat different notion of what information should be accessible in $Z$.
That is why different self-supervised methods can lead to different probing behavior, adaptation behavior, and domain suitability.

A good representation is therefore not simply one that minimizes a pretraining loss.
It is one whose induced hypothesis class
\[
\{g \circ \phi : g \in \mathcal{G}\}
\]
supports many downstream tasks with simple readout or modest adaptation.
This connects self-supervised learning back to the representation-learning picture of Section 4.1.
The pretext task is not the final goal.
It is an instrument for shaping a useful latent coordinate system.

\subsection*{What to remember}

\begin{itemize}
    \item Self-supervised learning constructs predictive targets from the data itself rather than from externally supplied labels.
    \item The three main objective families are masked prediction, autoregressive prediction, and contrastive learning.
    \item Augmentations, positives, negatives, and contexts are not minor implementation details; they determine what invariances and dependencies the representation is encouraged to learn.
    \item Contrastive learning has real theorem-level support through results such as the InfoNCE lower-bound interpretation, but many practical design choices remain heuristic.
    \item The reason self-supervision matters is not only that it avoids labels, but that it can induce reusable internal coordinates for downstream tasks.
\end{itemize}

\subsection*{Common confusion}

\begin{itemize}
    \item \textbf{``Self-supervised'' does not mean unsupervised in the sense of having no target at all.} There is a target, but it is generated from the example itself.
    \item \textbf{Low pretraining loss is not the same as good transfer.} A method can optimize its pretext objective well while learning the wrong invariances for downstream use.
    \item \textbf{Contrastive learning is not only geometry and not only information theory.} In practice it has both interpretations, and neither fully replaces the other.
    \item \textbf{Masked prediction and autoregression are related but not identical.} They expose different conditional structure and often produce different downstream behavior.
\end{itemize}

\subsection*{Exercises}

\begin{enumerate}
    \item Write the masked-prediction loss for a sequence in which exactly one token position is chosen uniformly at random for masking.
    \item Show that autoregressive next-token prediction is a maximum-likelihood objective for a factorized model of a sequence.
    \item In the one-negative contrastive example above, recompute the loss when $s(z,z^+) = 1.5$, $s(z,z^-) = 1.5$, and explain the result geometrically.
    \item Give one example of an augmentation that is likely valid for natural images and one that is likely invalid, and explain what invariance each would encourage.
    \item Compare masked language modeling and contrastive image learning as representation-learning mechanisms. What information is each objective trying to preserve, and what information is each willing to discard?
    \item Suppose a representation collapses so that $\phi(x)$ is the same vector for every $x$. Explain why this is bad for contrastive learning and for downstream linear probing.
    \item Derive from first principles why the softmax cross-entropy in contrastive learning rewards larger positive similarity than negative similarity.
    \item Consider a time series with latent periodic structure. Which of the three objective families in this section would you try first, and why?
\end{enumerate}

\subsection*{Notes and references}

Self-supervised learning has roots in several traditions rather than a single founding paper.
Predictive coding, language modeling, denoising autoencoding, masked prediction, and contrastive estimation all contributed pieces of the modern picture.
For denoising and representation learning, a classical reference is Vincent et al.
For language-model pretraining, early neural language-model and contextual-representation work is foundational.
For masked language modeling and bidirectional contextual pretraining, Devlin et al. is central.
For contrastive objectives and their information-theoretic interpretation, the Noise-Contrastive Estimation and InfoNCE lines are especially important.
For modern contrastive image learning, representative references include methods in the SimCLR and MoCo families.
A useful reading strategy is to separate three questions while reading the literature: what predictive task is defined, what invariances are encoded through the data pipeline, and what evidence is provided that the learned representation transfers.
